% $Header: /home/vedranm/bitbucket/beamer/solutions/conference-talks/conference-ornate-20min.en.tex,v 90e850259b8b 2007/01/28 20:48:30 tantau $

\documentclass{beamer} %[compress]
% alternatively add handout or draft option

%\mode<handout>%{\setbeamercolor{background canvas}{bg=black!5}}
%\usepackage{pgfpages}
%\pgfpagesuselayout{4 on 1}[a4paper,border shrink=5mm,landscape]
% For handout printing, 4 slides on 1 page

\mode<presentation>
{
  \usetheme{Pittsburgh} % plain template Pittsburgh
  \useoutertheme{infolines} % bottom (author, title, place)
  \useoutertheme[subsection=false]{miniframes} % top (sections, frames)
  \usecolortheme{beaver} %blue color whale
	\setbeamertemplate{frametitle}[default][right] % frame titles on the right
  \setbeamercovered{transparent}
  % or whatever (possibly just delete it)
}

\setbeamersize{text margin left=1cm, text margin right=1cm}
%\setbeamertemplate{navigation symbols}{}
\usepackage[english]{babel}
\usepackage{booktabs}
\usepackage{comment}
% or whatever

\usepackage[cp1250]{inputenc}
% or whatever
\usepackage{subfigure} %subfigures

\usepackage{times}
%\usepackage{breqn}
\usepackage{amsmath}
\usepackage[T1]{fontenc}
\def\sym#1{\ifmmode^{#1}\else\(^{#1}\)\fi} %shortcut for Stata tables
% Or whatever. Note that the encoding and the font should match. If T1
% does not look nice, try deleting the line with the fontenc.

%\logo{\includegraphics[height=0.5cm]{logo.jpg}} % logo on every page



\title[Lecture 4: Data Collection] % (optional, use only with long paper titles)
{Lecture 4: Data Collection}

\author[T.\ Havranek] % (optional, use only with lots of authors)
{ Tomas Havranek
}
% - Give the names in the same order as the appear in the paper.
% - Use the \inst{?} command only if the authors have different
%   affiliation.

\institute[Prague] % (optional, but mostly needed)
{
  	Charles University, CEPR, METRICS

  %\inst%
  %Economic Research Department\\
  %Czech National Bank
      %\inst{1}%
  %Research Department\\
  %Czech National Bank
  %\and
  %\inst{2}%
  %Institute of Economic Studies\\
  %Charles University, Prague
  
%\pgfdeclareimage[height=1cm]{logo}{logo.jpg} % logo only on title page
%\pgfuseimage{logo}
}% - Use the \inst command only if there are several affiliations.
% - Keep it simple, no one is interested in your street address.

\date[Meta in Econ and Finance] % (optional, should be abbreviation of conference name)
{Research Synthesis in Economics and Finance, \\ University of Canterbury\\[2ex]
  Term 1, 2025
}%\\[4mm]% - Either use conference name or its abbreviation.
% - Not really informative to the audience, more for people (including
%   yourself) who are reading the slides online


%\subject{Intertemporal Substitution}

\begin{comment}
\AtBeginSection[]
{
  \begin{frame}<beamer>{Outline}
    \tableofcontents[currentsection]%,currentsubsection]
  \end{frame}
}
\end{comment}

\begin{document}

\begin{frame}
  \titlepage
\end{frame}

\begin{comment}

\begin{frame}{Outline}
  \tableofcontents
  % You might wish to add the option [pausesections]
\end{frame}

\end{comment}

%\section{Summary}
%\subsection{}

%\section{EIS}
%\subsection{}


\begin{frame}{Effects must be quantitatively comparable}
What does it mean?

\bigskip

\begin{itemize}
	\item When we compare estimates reported in Study 1 and Study 2, it must be clear which one is larger
	\item It must be clear how much larger it is
	\item t-statistics are not useful measures of effect size
	\item Regression coefficients are not comparable in general if studies use different units
\end{itemize}

\end{frame}



\begin{frame}{Elasticities}

\[
\text{x-elasticity of } y: \, \varepsilon = \frac{\frac{\partial y}{y}}{\frac{\partial x}{x}}
\]

\[
\text{P-elasticity of Q: } \varepsilon = \frac{\frac{\partial Q}{Q}}{\frac{\partial P}{P}} \, \text{if continuous}
\]

\[
\varepsilon = \frac{\frac{Q_2 - Q_1}{Q_1}}{\frac{P_2 - P_1}{P_1}} = \frac{\text{\% change in quantity Q}}{\text{\% change in price P}} \, \text{if discrete}
\]

Elasticities can be recovered from regressing \textbf{log(Y) on log(X)}.


\end{frame}


\begin{frame}{Semi-elasticities, dollar values, std.dev.\ changes}
\begin{itemize}
	\item Semi-elasticity: effect of euro adoption on trade, effects of foreign investment on domestic productivity, effects of borders on trade
	\item Dollar values: statistical value of life, social cost of carbon
	\item Std.dev.\ changes: effect of class size on student test scores. What to do here?
	\item Class size is measured the same way in all studies (number of kids).
	\item Test scores are measured differently.
	\item Solution: By how many standard deviations do test scores improve if we decrease class size by 10 students?
\end{itemize}
\end{frame}


\begin{frame}{Standardized mean differences}

\[
\theta = \frac{\mu_1 - \mu_2}{\sigma},
\]
where $\mu_1$ is the mean for one population, $\mu_2$ is the mean for the other population, and $\sigma$ is a standard deviation based on either or both populations.

\bigskip
\begin{itemize}
	\item Very popular in psychology, not much used in economics. Why?
	\item Generally we need an experiment: treatment group, control group.
	\item In economics we often focus on observational research.
	\item For many questions the treatment and control are not binary classifications (we look at the intensity of treatment)
\end{itemize}


\end{frame}



\begin{frame}{Partial correlation coefficients}

Example: \textbf{tuition fees and university enrollment}. Ideally, we would need price elasticities of demand (PED), but many studies are not log on log.

\smallskip

 So we transform the estimates to partial correlations using t-statistics and degrees of freedom:

\bigskip

\[
\text{PCC(PED)}_{ij} = \frac{\text{T(PED)}_{ij}}{\sqrt{\text{T(PED)}_{ij}^2 + \text{DF(PED)}_{ij}}}.
\]



\end{frame}



\begin{frame}{PCC caveats}
\begin{itemize}
	\item Last resort: can be used almost always, but we lose a lot of information
	\item Difficult to interpret (some guidelines provided by Doucouliagos, 2011)
		\item Statistical problems for meta-analysis methods (but this is the case also for SMD)
	\item Include robustness checks: e.g.\ a subsample of elasticities in the case of tuition fees and university enrollment

\end{itemize}


\end{frame}



\begin{frame}{Collect all estimates}
\begin{figure}
	\centering
		\includegraphics[width=0.9\textwidth]{pref.png}
\end{figure}
\end{frame}


\begin{frame}{Standard errors}
\begin{itemize}
	\item Absolutely crucial for meta-analysis
	\item If not reported directly, can be computed from t-statistics or p-values
	\item Last resort: approximation based on within-study variation (Matousek et al.\, 2022)
	\item Interactions, transformations: delta method
\end{itemize}
\end{frame}

\begin{frame}{Delta method}
Based on Taylor series expansion. Basic formula:

\[
\text{Var}(Y) = \text{Var}(f(X)) \approx [f'(\mu_X)]^2 \text{Var}(X).
\]

% So, if $Y$ is any function of a random variable $X$, we need only calculate the variance of $X$ and the first derivative of the function to approximate the variance of $Y$.

\bigskip

\underline{\text{Example:}} \, \text{Suppose } $Y = X^2$. \, \text{Then } $f(x) = x^2$ \, \text{and } $f'(x) = 2x$, \, \text{so that:}
\[
\text{Var}(Y) \approx [2\mu_X]^2 \text{Var}(X) = 4\mu_X^2\sigma_X^2.
\]

\underline{\text{Example:}} \, \text{Suppose } $Y = \frac{1}{X}$. \, \text{Then } $f(x) = \frac{1}{x}$ \, \text{and } $f'(x) = -\frac{1}{x^2}$, \, \text{so that:}
\[
\text{Var}(Y) \approx \left[\frac{-1}{\mu_X^2}\right]^2 \text{Var}(X) = \frac{\sigma_X^2}{\mu_X^4}.
\]



\end{frame}


\begin{frame}{Graphical estimates, measurement error}
\begin{figure}
	\centering
		\includegraphics[width=1\textwidth]{measure.png}
\end{figure}
\end{frame}


\begin{frame}{Outliers}
\begin{figure}
	\centering
		\includegraphics[width=0.9\textwidth]{outlier.png}
\end{figure}
\end{frame}


\begin{frame}{What to do with outliers?}
\begin{itemize}
	\item No consensus
	\item Approach 1: omit very large effects (and standard errors), e.g.\ more than 3 standard deviations away from the mean
	\item Approach 2: run robust meta-regression or use tools for outlier detection in regression
	\item Approach 3: use winsorizing (my personal preference)
	\item Always report robustness checks
\end{itemize}
\end{frame}


\begin{frame}{Winsorization}
\begin{figure}
	\centering
		\includegraphics[width=0.65\textwidth]{winsor.png}
\end{figure}
\end{frame}



\begin{frame}{Reflecting context: the effect of beauty on success}
\begin{figure}
	\centering
		\includegraphics[width=1\textwidth]{data.png}
\end{figure}
\end{frame}

\begin{frame}{Measurement}
\begin{figure}
	\centering
		\includegraphics[width=1\textwidth]{measure1.png}
\end{figure}
\end{frame}


\begin{frame}{Data, estimation, publication}
\begin{figure}
	\centering
		\includegraphics[width=1\textwidth]{est.png}
\end{figure}
\end{frame}


\begin{frame}{Artificial intelligence and context}
\begin{itemize}
	\item Can be useful for identifying key aspects in which studies vary
	\item Do not trust AI replies too much, but you can use it at various stages to generate ideas about potential explanatory variables
	\item You can use it to check whether you missed an important variable
\end{itemize}
\end{frame}


\begin{frame}{Value added}
\begin{itemize}
	\item Meta-analysis can be used to test hypotheses that primary studies cannot investigate
	\item Example: the effect of daylight saving time on energy consumption
	\item Individual studies have data for individual countries (or provinces)
	\item In meta-analysis, we have data for dozens of countries
	\item We can investigate the sources of cross-country differences
\end{itemize}
\end{frame}


\begin{frame}{DST savings and latitude}
\begin{figure}
	\centering
		\includegraphics[width=1\textwidth]{bma_uip.pdf}
\end{figure}
\end{frame}


\begin{frame}{House prices and monetary policy}
\begin{figure}
	\centering
		\includegraphics[width=0.8\textwidth]{bma.pdf}
\end{figure}
\end{frame}



\end{document}


